\scanpage{12}
\inserted{à scalaire près} \deleted{fonction quadratique sur}\inserted{poly. de degré $\leq2$} $\unsure{V|F}$ invariante par $\mathbb D_\infty$ et \deleted{\unsure{non-dégén}}nulle sur les sommets du polygône $P$, savoir
\[
xy-1
\]
[ou encore, $\exists!$ forme quadratique sur $V$ invariante par $\mathbb D_\infty$ égale à $1$ sur les sommets du polygône $P$, savoir $xy\ldots$ ].

\uline{NB} Il y a ``en général'' une et une seule cônique qui passe par cinq points. (\uline{NB} cinq équations lin. \unsure{strictes|scalaires|homogènes} sur $6$ coeff. homogènes...), donc \deleted{on} s'attend à ce \deleted{que}, lorsque $P$ a au moins cinq pts fibre par fibre (i.e. $\alpha\ne0,-1$ fibre par fibre pour exclure triangles et carrés, donc au total
\[
\alpha(\alpha^2-4)(\alpha+1)\quad\text{inversible}
\]
qu'il existe \deleted{\unsure{$\exists$}}\inserted{une unique} conique qui passe par les sommets de $P$. Écrivant la cônique sous forme non homogène comme
\[
Q(x,y)=ax^2+\deleted{\unsure{2}}bxy+cy^2+ux+vy+\lambda
\]
et prenons pour $(x,y)$ les cinq points $({\unsure{\zeta|\xi}}^i,\zeta^{-i})$, $-2\leq i\leq+2$, on trouve
% archive sheet number: 11
